\chapter{Inflation, structure et perte des conditions initiales}

\section{{\proseunicodefont Ω} et les indices en faveur de la platitude}

Une question d’étudiant sur le grand \(\Omega\) ramène la discussion à l’équation de Friedmann-Robertson-Walker. À l’époque actuelle, bien après la domination du rayonnement, la loi d’expansion s’écrit avec un terme d’énergie du vide et un terme de matière ; le rayonnement est omis parce qu’il est désormais très faible :
\begin{equation}
H^2=\left(\frac{\dot a}{a}\right)^2=\frac{8\pi G}{3}\left(\rho_0+\rho_m\right)-\frac{k}{a^2}.
\end{equation}
Le terme du vide ne se dilue pas. Le terme de matière, lui, se dilue :
\begin{equation}
\rho_m=\frac{M}{a^3}.
\end{equation}
\(\Omega\) représente la fraction de l’équation d’expansion associée aux différentes composantes.

Divisons par \(H^2\). Les termes normalisés ont alors une somme égale à un. Avec une notation régularisée,
\begin{equation}
\begin{aligned}
\Omega_0&=\frac{8\pi G\rho_0}{3H^2},\\
\Omega_m&=\frac{8\pi G\rho_m}{3H^2},\\
\Omega_k&=-\frac{k}{a^2H^2},
\end{aligned}
\end{equation}
et par conséquent
\begin{equation}
\Omega_0+\Omega_m+\Omega_k=1.
\end{equation}
\editorialnote{L’enregistrement audio et le tableau sont difficiles à interpréter dans ce passage, et l’orateur emploie plusieurs noms pour désigner la contribution du vide. La notation présentée ici régularise le raisonnement après division par \(H^2\).}
Si l’espace est plat, \(k=0\), donc
\begin{equation}
\Omega_0+\Omega_m=1.
\end{equation}
Si la somme des contributions de la matière et du vide est supérieure à un, la courbure est positive. Si elle est inférieure à un, la courbure est négative.

La logique observationnelle est alors simple. On mesure la matière ordinaire et la matière noire. On mesure l’expansion de Hubble. On mesure géométriquement la courbure à très grande échelle. On constate aussi que l’Univers est en accélération. Le signe du terme du vide n’est pas déterminé par \(H\) à lui seul, mais par la tendance de \(\dot a/a\) vers une constante, ce qui signifie que le facteur d’échelle tend vers la branche exponentielle,
\begin{equation}
a(t)\propto e^{Ht}.
\end{equation}
Ces différentes mesures se recoupent et indiquent que l’Univers est plat à quelques pour cent près.

\begin{classroomqa}
\lecturetimestamp{00:10:50}
\audiencequestion{La matière comprend-elle la matière noire ?}
\lecturerresponse{Oui. Dans la discussion présente, le terme de matière comprend la matière noire. On pourrait aussi ajouter le rayonnement, mais sa fraction est aujourd’hui si faible qu’on peut la négliger.}
\end{classroomqa}

\begin{classroomqa}
\lecturetimestamp{00:11:11}
\audiencequestion{Existe-t-il une mesure de l’énergie noire autre que la constante de Hubble ?}
\lecturerresponse{Ce qui importe n’est pas la valeur de \(H\) à elle seule, mais la façon dont l’expansion tend vers une loi exponentielle. Avec la mesure de la matière et la mesure géométrique de la platitude, cela permet d’inférer le terme du vide.}
\end{classroomqa}

\section{De la platitude à la structure}

Pourquoi l’Univers est-il si plat ? L’hypothèse avancée est celle de l’inflation. Plat signifie très grand, très lisse et très homogène sur des échelles de milliards, en réalité de dizaines de milliards, d’années-lumière. L’expansion exponentielle peut étirer une région jusqu’à lui donner cet aspect.

Mais on peut pousser trop loin une bonne chose. Un Univers parfaitement lisse le reste. La gravité ne crée pas d’irrégularités à partir d’une symétrie exacte. Si des structures existent aujourd’hui, elles ont dû apparaître plus tôt à partir d’une faible inhomogénéité. En remontant dans le temps à partir de l’Univers actuel, on trouve qu’à l’époque de la transparence, l’inhomogénéité était de l’ordre d’une partie sur \(10^5\) :
\begin{equation}
\frac{\delta\rho}{\rho}\sim 10^{-5}.
\end{equation}
L’Univers était extrêmement lisse, mais pas parfaitement homogène.

À partir de là, la gravité agit contrairement aux forces ordinaires dans une pièce remplie d’air. Une légère surdensité attire la matière vers elle. La région voisine s’appauvrit ; la sous-densité devient donc encore moins dense, tandis que la surdensité s’accentue. La gravité renforce le contraste au lieu de l’effacer.

La question naturelle est alors de savoir pourquoi cela ne s’est pas produit plus tôt. La réponse est la pression. Plus tôt, la pression de rayonnement était élevée. En utilisant
\begin{equation}
p=w\rho,
\end{equation}
les cas familiers sont
\begin{equation}
w_{\rm rad}=\frac13,\qquad
w_m=0,\qquad
w_{\rm vac}=-1.
\end{equation}
Une pression positive homogénéise. Elle repousse les perturbations et lisse le milieu. À grande échelle, la matière étant traitée comme dépourvue de pression, elle n’offre aucun soutien de ce type. Tant que la pression est présente, elle contrecarre la gravité. Quand elle devient négligeable, la gravité l’emporte.

La matière noire joue un rôle essentiel dans l’ordre des événements. Les premiers amas à s’effondrer sont des halos de matière noire. La matière baryonique ordinaire tombe ensuite dans ces halos et, comme elle peut dissiper de l’énergie, se concentre plus près du centre. La matière noire ne dissipe pas de la même manière et reste donc répartie dans le halo. Le début de cette contraction précède légèrement l’époque de la transparence, d’un facteur d’environ dix en temps, ce qui explique pourquoi le fond diffus cosmologique n’en révèle pas directement le début.

Une brève digression précise l’histoire chimique ultérieure. Presque tous les éléments plus lourds que l’hélium sont issus de l’évolution stellaire, et non de la cosmologie de l’Univers primordial. L’Univers primordial produit essentiellement de l’hydrogène, de l’hélium et de faibles traces de quelques noyaux légers.

Ce que l’on observe dans le ciel n’est pas directement la densité de masse, mais l’anisotropie de température du fond diffus cosmologique. Les fluctuations de densité y contribuent, tout comme les fluctuations de vitesse. Les régions surdenses paraissent légèrement plus froides, car les photons perdent davantage d’énergie en remontant des puits gravitationnels plus profonds ; les décalages Doppler dus au mouvement du gaz s’y superposent. L’analyse numérique complète est complexe. L’énoncé robuste est plus simple : au moment où l’Univers est devenu transparent, la densité d’énergie variait d’environ une partie sur \(10^5\).

\begin{classroomqa}
\lecturetimestamp{00:20:29}
\audiencequestion{Pourquoi la gravité n’a-t-elle pas fait croître les amas plus tôt ?}
\lecturerresponse{Parce que l’Univers était auparavant dominé par le rayonnement, dont la pression positive était suffisante pour homogénéiser le milieu. La gravité n’a commencé à dominer la croissance des amas que lorsque le rayonnement est devenu négligeable devant la matière.}
\end{classroomqa}

\begin{classroomqa}
\lecturetimestamp{00:24:35}
\audiencequestion{L’effondrement ne nécessite-t-il pas un mécanisme de dissipation ?}
\lecturerresponse{La matière baryonique ordinaire a besoin de dissiper de l’énergie pour se déposer vers le centre. Mais la matière noire peut commencer à former des halos sans ce type de dissipation collisionnelle. Il existe aussi une forme d’amortissement cosmologique liée à l’expansion elle-même, qui jouera un rôle plus tard.}
\end{classroomqa}

\begin{classroomqa}
\lecturetimestamp{00:30:34}
\audiencequestion{Y avait-il aussi des fluctuations de vitesse, permettant aux galaxies d’acquérir ensuite un moment cinétique ?}
\lecturerresponse{Oui. Le motif primordial comportait des fluctuations de densité et des fluctuations de vitesse. Ce que l’on observe dans le ciel est un motif de température produit à la fois par les variations de densité et par le mouvement du gaz.}
\end{classroomqa}

\section{L’inflaton et le plateau}

Le mécanisme présenté maintenant repose sur un champ scalaire \(\phi\). Sa densité d’énergie comporte trois termes :
\begin{equation}
\rho_\phi=\frac12\dot\phi^2+\frac12(\nabla\phi)^2+V(\phi).
\end{equation}
La variation temporelle contribue. La variation spatiale contribue également. Un champ scalaire peut aussi posséder une véritable énergie potentielle \(V(\phi)\). Cette contribution n’a rien d’exotique : le champ de Higgs en est l’exemple familier. Ce qui importe ici, c’est la forme du potentiel.

Le potentiel doit être assez plat sur une certaine plage, avec une légère pente descendante ; puis, à une valeur critique, le fond s’effondre. Il ne tombe pas exactement à zéro, mais à une valeur très faible. En unités de Planck, l’énergie actuelle du vide est de l’ordre de
\begin{equation}
10^{-123}.
\end{equation}
Sur le plateau, elle était bien plus élevée, tout en restant très inférieure à l’échelle de Planck.

\begin{figure}[t]
\centering
\includegraphics[width=0.78\textwidth]{lecture_08_figure_03.png}
\caption{Potentiel qualitatif de l’inflaton : large plateau, chute abrupte et minimum bas.}
\lectureframe{00:43:23}
\end{figure}

Ce graphique ne représente pas un lieu dans l’espace. Il représente une valeur du champ. Dire que l’Univers a commencé en haut du plateau signifie seulement que, dans une très grande région, la valeur du champ se trouvait là-haut. Peu importe si \(\phi\) variait quelque peu d’un endroit à l’autre au départ, car l’inflation étire énormément ces variations. Après une expansion suffisante, chaque région locale paraît presque uniforme. \editorialnote{Les axes ne sont pas lisibles sur l’image du tableau, mais l’explication orale indique clairement que l’axe vertical représente l’énergie potentielle \(V(\phi)\) et l’axe horizontal la valeur du champ \(\phi\).}

Tant que le champ se trouve sur la partie large et de faible pente du potentiel, son énergie potentielle se comporte comme l’énergie du vide. Si \(\phi\) était parfaitement figé, il se comporterait exactement comme l’énergie du vide. En omettant les facteurs numériques non indiqués, la relation donnée dans le cours est simplement
\begin{equation}
H\propto \sqrt{V(\phi)},
\end{equation}
ce qui donne la loi presque exponentielle
\begin{equation}
a(t)\propto e^{Ht}.
\end{equation}
Presque tout ce que l’on peut observer aujourd’hui se produit après que le champ a franchi le bord. C’est au cours de cette chute que la grande énergie, semblable à celle du vide sur le plateau, se transforme en rayonnement, en matière ordinaire, en matière noire et en la faible énergie du vide résiduelle observée aujourd’hui.

\begin{classroomqa}
\lecturetimestamp{00:44:07}
\audiencequestion{Sur quel espace-temps le champ \(\phi\) est-il défini, et utilise-t-on toujours les équations d’Einstein ?}
\lecturerresponse{Oui. \(\phi\) est un champ défini sur l’espace-temps, et on utilise toujours la même équation de Friedmann. Dans ce régime, la densité d’énergie au membre de droite est principalement fournie par \(V(\phi)\), à condition que le champ varie assez lentement.}
\end{classroomqa}

\begin{classroomqa}
\lecturetimestamp{00:47:04}
\audiencequestion{Comment l’énergie de l’inflaton, semblable à celle du vide, peut-elle se transformer en matière et en rayonnement ?}
\lecturerresponse{Quand le champ descend jusqu’au fond et oscille, ces oscillations sont les quanta du champ. Si ceux-ci sont instables, ils peuvent se désintégrer en photons, électrons, positrons et autres particules familières.}
\end{classroomqa}

\section{Roulement lent et frottement de Hubble}

Une fois que le champ a franchi le bord, pourquoi ne continuerait-il pas simplement à osciller indéfiniment ? Pourquoi ne ferait-il pas des allers-retours, voire ne remonterait-il pas ? La réponse donnée ici est un mécanisme d’amortissement cosmologique. Il ne s’agit pas d’un véritable frottement microscopique, mais le terme correspondant dans l’équation du mouvement agit comme tel.

Négligeons les gradients spatiaux et considérons le champ comme homogène. Le terme cinétique et le terme de potentiel se comportent alors comme un degré de liberté unique avec une dynamique \(T-V\), mais comme il s’agit de densités d’énergie, ils doivent être multipliés par le volume physique \(a^3\). On obtient un lagrangien effectif homogène de la forme
\begin{equation}
L_{\mathrm{hom}}=a^3(t)\left[\frac12\dot\phi^2-V(\phi)\right].
\end{equation}
\editorialnote{La ligne supérieure du tableau est partiellement coupée. Ce lagrangien homogène n’est reconstitué que dans la mesure nécessaire pour retrouver l’équation du mouvement fiable inscrite juste en dessous.}

L’équation lisible au tableau est
\begin{equation}
\frac{d}{dt}\left(a^3\dot\phi\right)=-a^3\frac{\partial V}{\partial\phi}.
\end{equation}
En développant la dérivée temporelle, on obtient
\begin{equation}
a^3\ddot\phi+3a^2\dot a\,\dot\phi=-a^3\frac{\partial V}{\partial\phi},
\end{equation}
et en divisant par \(a^3\), il vient
\begin{equation}
\ddot\phi+3H\dot\phi=-\frac{\partial V}{\partial\phi}.
\end{equation}
Le signe moins est corrigé explicitement dans la démonstration orale : la force est l’opposé de la dérivée de l’énergie potentielle.

\begin{figure}[t]
\centering
\includegraphics[width=0.78\textwidth]{lecture_08_figure_05.png}
\caption{Équation du mouvement homogène de l’inflaton, avec le terme de frottement de Hubble \(3H\dot\phi\).}
\lectureframe{00:57:10}
\end{figure}

Le terme \(3H\dot\phi\) est ce qu’on appelle le frottement de Hubble. Il provient entièrement de la dérivation du facteur \(a^3\). Du point de vue de la dynamique, il ressemble toutefois au terme de traînée qui s’exerce sur un rocher tombant dans l’eau. Il y a un terme de force, un terme d’accélération et un terme linéaire en la vitesse. Le champ atteint rapidement une vitesse limite, puis évolue lentement le long du plateau. C’est le régime de roulement lent. Tant que la pente est faible et que le champ se déplace lentement, l’énergie potentielle ne varie que lentement et l’expansion reste presque exponentielle.

Près du minimum, le champ oscille. Ces oscillations sont les quanta de l’inflaton. S’ils sont instables, ils se désintègrent, ce qui transforme l’énergie du plateau en matière chaude et en rayonnement dans l’Univers ultérieur.

\begin{classroomqa}
\lecturetimestamp{00:50:48}
\audiencequestion{Pourquoi le champ n’oscille-t-il pas indéfiniment, et qu’est-ce qui produit l’amortissement ?}
\lecturerresponse{L’expansion elle-même produit un terme semblable à un frottement. Le facteur \(a^3\) dans le lagrangien homogène engendre le terme \(3H\dot\phi\), qui agit comme une traînée visqueuse.}
\end{classroomqa}

\section{Fluctuations quantiques et sortie de l’horizon}

Un problème apparent se pose alors. Si l’inflation rend l’Univers extraordinairement lisse, pourquoi n’efface-t-elle pas les germes nécessaires aux galaxies ? La réponse est la mécanique quantique. Les fluctuations qui subsistent ne sont pas les vestiges d’un lissage imparfait : elles sont produites à nouveau.

Oublions un instant l’expansion. Même si un champ était presque parfaitement homogène, la mécanique quantique ne permettrait pas de fixer exactement à la fois \(\phi\) et \(\dot\phi\). L’état fondamental reste agité. C’est la même raison pour laquelle un oscillateur harmonique ordinaire possède une énergie de point zéro. Dans les circonstances habituelles, ces fluctuations sont difficiles à observer parce qu’elles ont une courte longueur d’onde et oscillent rapidement ; les observations en font donc la moyenne.

L’expansion change la situation. En coordonnées comobiles, le mode du champ conserve à peu près la même longueur d’onde coordonnée, tandis que sa longueur d’onde physique croît avec le facteur d’échelle :
\begin{equation}
\lambda_{\rm phys}=a(t)\,\Delta x.
\end{equation}
Pendant l’inflation, l’échelle physique de Hubble est approximativement constante,
\begin{equation}
d_H\sim H^{-1},
\end{equation}
si bien que l’horizon comobile correspondant se contracte comme
\begin{equation}
d_{H,\mathrm{coord}}\sim \frac{1}{aH}.
\end{equation}

L’analogie de la corde permet de comprendre le gel des modes. Une onde oscille parce que chaque segment de la corde est tiré par ses voisins. Mais si sa longueur d’onde devient aussi grande que l’horizon, les voisins concernés se trouvent au-delà du contact causal. La force de rappel disparaît. Dès qu’un mode dépasse l’horizon, il se fige.

Cela se produit encore et encore. Les grandes longueurs d’onde se figent d’abord. Les plus courtes se figent plus tard. Même un mode dont la longueur d’onde est initialement minuscule finit par être étiré jusqu’à franchir l’horizon. L’inflation n’efface donc pas simplement la structure : elle continue à produire des variations du champ qui se figent.

Le processus ne s’arrête pas non plus parce que l’inflation dure longtemps. Les anciens modes figés sont étirés jusqu’à des échelles immenses, mais de nouveaux modes continuent de se figer. Le motif accumulé devient statistiquement stationnaire au sens pertinent en cosmologie. C’est l’origine du spectre approximativement invariant d’échelle : le même processus de franchissement de l’horizon se répète à des échelles logarithmiques successives et produit à peu près la même amplitude d’une bande de longueurs d’onde à la suivante.

Après que le champ a quitté le plateau, l’horizon recommence à croître. Les modes qui se trouvaient au-delà de l’horizon inflationnaire temporaire y rentrent par la suite. Selon la formulation habituelle, ils sortent de l’horizon pendant l’inflation, puis y rentrent après celle-ci. Ils ont alors été étirés jusqu’à des échelles macroscopiques et, lorsqu’ils recommencent à osciller, ils le font sous forme de structures de très grande taille.

\begin{classroomqa}
\lecturetimestamp{01:12:02}
\audiencequestion{À quoi ressemble l’horizon sur le diagramme d’espace-temps ?}
\lecturerresponse{Pendant l’expansion exponentielle, la distance physique à l’horizon reste à peu près égale à \(c/H\), mais, sur un diagramme comobile, cette distance physique fixe paraît de plus en plus petite, car l’échelle des coordonnées est étirée. En coordonnées comobiles, l’horizon se contracte comme \(1/(aH)\).}
\end{classroomqa}

\begin{classroomqa}
\lecturetimestamp{01:15:59}
\audiencequestion{Que devient l’amplitude lorsque les fluctuations sont étirées, et existe-t-il un régime stationnaire ?}
\lecturerresponse{Les modes individuels se diluent à mesure qu’ils sont étirés, mais de nouveaux modes continuent de se figer. Il en résulte un motif statistique stationnaire, approximativement invariant d’échelle, plutôt qu’une disparition complète des fluctuations.}
\end{classroomqa}

\section{Des variations gelées du champ aux variations de densité}

Une fois que le champ a acquis de faibles variations d’un endroit à l’autre, celles-ci se transforment en différences de densité lorsque les différentes régions franchissent le bord à des instants légèrement différents. Une région un peu plus bas sur la pente atteint la falaise en premier. Si son énergie du vide se transforme en rayonnement, celui-ci se dilue avec l’expansion. Une région voisine qui se trouve encore sur le plateau reste dominée par l’énergie du vide et ne se dilue pas. Lorsque la seconde région chute à son tour, la première a déjà perdu une partie de son énergie par décalage vers le rouge.

C’est pourquoi un roulement plus lent produit des contrastes de densité finaux plus importants. Si le champ avance très lentement le long du plateau, une faible différence de valeur du champ correspond à un décalage temporel plus grand entre le moment où une région atteint la falaise et celui où sa voisine l’atteint. La région qui chute en premier a davantage de temps pour se diluer avant que l’autre ne convertisse son énergie. Plus la pente est faible, plus la vitesse limite est basse et plus la variation finale de densité est grande. Un plateau plus élevé tend également à produire des fluctuations plus importantes. La largeur du plateau détermine principalement la durée de l’inflation.

Seule une combinaison de la hauteur et de la pente est étroitement contrainte par les données sur les fluctuations. Si l’échelle d’énergie est assez élevée, une contribution d’ondes gravitationnelles devrait aussi être observable. Les limites numériques ne sont données qu’à titre d’ordre de grandeur. En gros, si la densité d’énergie sur le plateau dépassait largement environ \(10^{-14}\) en unités de Planck, le fond diffus cosmologique présenterait déjà trop de structures dues aux ondes gravitationnelles ; si elle était nettement inférieure à environ \(10^{-16}\), il y aurait peu de chances réalistes de détecter un tel signal. \editorialnote{L’orateur précise explicitement que ces nombres sont approximatifs et qu’il ne se souvient pas de leurs valeurs exactes. Ils ne sont conservés ici qu’à titre d’ordres de grandeur.}

Les unités sont expliquées selon la méthode habituelle des unités naturelles. La longueur d’onde de Compton d’une particule de masse \(m\) est
\begin{equation}
\lambda_C=\frac{\hbar}{mc}.
\end{equation}
Si \(c=\hbar=1\), la masse et l’inverse d’une longueur ont les mêmes unités. La densité d’énergie étant l’énergie par unité de volume, elle s’exprime en unités de masse à la puissance quatre :
\begin{equation}
[\rho]=(\mathrm{mass})^4.
\end{equation}

En un point donné, l’accumulation de modes gelés ressemble à une marche aléatoire. Différents points de l’espace suivent des marches aléatoires différentes, ce qui transforme précisément le bruit quantique en variations spatiales.

\begin{classroomqa}
\lecturetimestamp{01:33:04}
\audiencequestion{\(\phi\) suit-il une marche aléatoire, et cela rend-il sans importance le point de départ antérieur ?}
\lecturerresponse{En un point fixe, les modes successifs qui se figent s’ajoutent avec des phases aléatoires ; \(\phi\) suit donc bien une marche aléatoire. Ce succès même pose aussi problème : l’inflation remplace sans cesse les informations antérieures par des fluctuations ultérieures, si bien qu’il devient difficile de retrouver ce qui s’est passé avant l’inflation.}
\end{classroomqa}

\section{Ce que l’inflation dissimule}

La théorie réussit dans un sens bien précis. Elle relie des fluctuations sur une vaste gamme d’échelles à un mécanisme unique. Mais ce même mécanisme efface ce qui l’a précédé. Les irrégularités antérieures sont étirées et les fluctuations quantiques produites plus tard sont celles qui restent pertinentes. Voilà pourquoi l’inflation est à la fois séduisante et frustrante. Elle explique l’homogénéité et le spectre des fluctuations, mais efface aussi une grande partie de l’histoire antérieure.

Cela suscite un malaise quant au point de départ. Une densité d’énergie très élevée ne pose pas problème en soi. Le problème est que l’énergie du plateau ne semble être naturellement ni de l’ordre de l’échelle de Planck ni extrêmement faible. Elle est élevée en unités ordinaires, mais reste faible en unités de Planck. Elle ne paraît pas être un point de départ évident pour une théorie fondamentale.

Dans cette représentation simple, il n’y a pas non plus de \(t=0\) au sens littéral. Dès que l’inflation a duré trop longtemps, presque toute information sur une étape antérieure a disparu. Si le nombre de plis d’expansion n’était que légèrement supérieur au minimum nécessaire, il subsisterait peut-être une trace de l’étape précédente ; en général, toutefois, l’inflation efface le passé.

La même réserve s’applique au traitement de l’espace de de Sitter. La loi exponentielle pure décrit correctement la branche inflationnaire aux temps tardifs, mais il vaut mieux décrire la solution exacte de de Sitter par
\begin{equation}
a(t)\propto \cosh(Ht).
\end{equation}
Cette solution se contracte jusqu’à une taille minimale, puis se dilate. La question de savoir s’il faut prendre au pied de la lettre la branche contractante en cosmologie reste ouverte.

\begin{classroomqa}
\lecturetimestamp{01:37:31}
\audiencequestion{Se pourrait-il qu’il n’y ait pas eu de commencement et que le comportement exponentiel se soit poursuivi indéfiniment vers le passé ?}
\lecturerresponse{Pas dans cette représentation simple. L’espace de de Sitter exact ne correspond pas à une expansion exponentielle éternelle à tout instant ; il ressemble plutôt à \(\cosh(Ht)\), avec une branche contractante et une branche en expansion. La question de savoir s’il faut prendre cette branche antérieure au pied de la lettre reste ouverte.}
\end{classroomqa}

\section{Ajustement fin, courbure et contraintes anthropiques}

En première approximation, les paramètres importants sont la hauteur et la pente du plateau, ainsi que sa largeur \(\Delta\phi\). La pente importe parce qu’elle fixe la vitesse limite et donc l’amplitude des fluctuations. La largeur importe parce qu’elle contribue à déterminer la durée de l’inflation. Selon un préjugé courant, le déplacement d’un champ fondamental ne devrait pas être beaucoup plus grand que la masse de Planck, et la densité d’énergie ne devrait pas dépasser l’ordre de l’unité en unités de Planck. Avec ces limites approximatives, une inflation suffisante n’est pas possible dans tout l’espace des paramètres. L’estimation citée ici correspond à un ajustement de quelques pour cent, environ \(2\%\), pour obtenir une pente assez faible et environ \(60\) à \(62\) plis d’expansion.

\begin{classroomqa}
\lecturetimestamp{01:46:02}
\audiencequestion{Peut-on comprendre l’inflation en termes d’espace des paramètres, et quel est le degré d’ajustement fin requis ?}
\lecturerresponse{Oui. En première approximation, on considère la hauteur, la pente et la largeur du plateau. Le cours donne une estimation approximative : obtenir suffisamment de plis d’expansion nécessite un ajustement de quelques pour cent, environ \(2\%\).}
\end{classroomqa}

Cet ajustement modéré est ensuite comparé à l’énigme bien plus aiguë de la constante cosmologique actuelle, de l’ordre de \(10^{-123}\). La discussion aborde les questions anthropiques, car ces deux problèmes peuvent être reformulés en demandant quelles conditions sont nécessaires à la formation des structures.

Dans le cadre newtonien, une constante cosmologique positive agit comme une répulsion universelle proportionnelle à la distance. Si elle était beaucoup plus grande, elle n’aurait pas besoin de disloquer les galaxies aujourd’hui pour avoir des conséquences fatales. Il suffirait qu’elle soit assez grande pour empêcher les fluctuations primordiales \(10^{-5}\) de se condenser en galaxies. L’exemple donné est qu’une valeur d’environ deux ordres de grandeur supérieure à la valeur observée suffirait déjà à empêcher cette formation. C’est le point associé à l’argument de Weinberg : la constante cosmologique n’a pas besoin d’être exactement nulle ; il suffit qu’elle soit assez faible pour permettre la formation de structures.

Une courbure négative agit dans le même sens. Dans le langage newtonien des cours précédents, \(k<0\) signifie que la vitesse de récession dépasse la vitesse de libération. Une poussée vers l’extérieur trop forte empêche elle aussi les faibles surdensités primordiales de devenir des galaxies. L’inflation dilue cette courbure. La limite inférieure citée est d’environ \(59.5\) plis d’expansion. Une valeur proche de \(62\) est alors considérée comme représentative d’une histoire inflationnaire suffisamment au-dessus du seuil.

Cela modifie la question de probabilité. Si l’on considère le volume brut de l’espace des paramètres qui permet une inflation suffisante, la réponse peut n’être que de quelques pour cent. Si l’on conditionne sur le fait de se trouver dans un Univers où des galaxies se forment et où des observateurs existent pour poser la question, la probabilité conditionnelle peut devenir bien plus grande. C’est sur ce point que s’achève la discussion : sans solution définitive, mais avec les problèmes d’ajustement fin, leur réinterprétation anthropique et le constat que les grands débats demeurent ouverts.
